%===========================================================================
% PART 2: SECTIONS C & D
%===========================================================================

\documentclass[aspectratio=169]{beamer}
\usetheme{Madrid}
\usecolortheme{default}
\usepackage{amsmath, amssymb, amsfonts}
\usepackage{graphicx}
\usepackage{booktabs}
\usepackage{array}
\usepackage{multirow}
\usepackage{bm}
\usepackage{xcolor}
\usepackage{tikz}
\usepackage{enumitem}
\usepackage{caption}
\usepackage{subcaption}
\usetikzlibrary{shapes,arrows,positioning,calc}

% Define colors
\definecolor{navyblue}{RGB}{0,53,102}
\definecolor{crimson}{RGB}{165,42,42}
\definecolor{darkgreen}{RGB}{0,100,0}
\definecolor{gold}{RGB}{184,134,11}
\definecolor{lightgray}{RGB}{240,240,240}
\definecolor{darkgray}{RGB}{100,100,100}

% Title formatting
\setbeamercolor{title}{fg=navyblue}
\setbeamercolor{frametitle}{fg=navyblue}
\setbeamercolor{structure}{fg=navyblue}
\setbeamercolor{block title}{bg=navyblue!20, fg=navyblue}
\setbeamercolor{alertblock title}{bg=crimson!20, fg=crimson}
\setbeamercolor{exampleblock title}{bg=darkgreen!20, fg=darkgreen}

% Custom commands - FIXED: Changed \note to \mynote
\newcommand{\question}[1]{{\large\bfseries\color{crimson}#1}}
\newcommand{\subquestion}[1]{{\bfseries\color{darkgreen}#1}}
\newcommand{\concept}[1]{\textcolor{navyblue}{\textbf{#1}}}
\newcommand{\formula}[1]{\textcolor{gold}{\textbf{#1}}}
\newcommand{\exampletext}[1]{\textit{\textcolor{darkgreen}{#1}}}
\newcommand{\highlight}[1]{\textcolor{crimson}{\textbf{#1}}}
\newcommand{\mynote}[1]{\textcolor{darkgray}{\textit{#1}}}

% Title Page
\title[Synoptic Examination]{Advanced Financial Theory \\[0.3cm]
	\large A Comprehensive Examination (Part 2)}
\subtitle{Sections C, D, and Conclusion}
\author{Shivgan Joshi}
\date{July 8, 2026}

\begin{document}
	
	%===========================================================================
	% CONTINUE FROM PART 1
	%===========================================================================
	\begin{frame}
		\titlepage
		\begin{center}
			\small \textit{``The capital structure decision is one of the most important and complex decisions a firm can make.''}
		\end{center}
	\end{frame}
	
	%===========================================================================
	% SECTION C: CAPITAL STRUCTURE
	%===========================================================================
	\section{Section C: Capital Structure and Strategic Finance}
	
	\begin{frame}{Section C: Capital Structure and Strategic Finance}
		\begin{center}
			\Huge\textcolor{navyblue}{\textbf{Section C}}
			\vspace{0.5cm}
			
			\huge Capital Structure and \\
			Strategic Finance
			\vspace{0.3cm}
			
			\large\textit{MM Propositions, Trade-Off Theory, and Pecking Order}
			\vspace{0.5cm}
			
			\begin{quotation}
				\small ``The capital structure decision is one of the most important and complex decisions a firm can make. The right capital structure can enhance value; the wrong one can destroy it.''
			\end{quotation}
		\end{center}
	\end{frame}
	
	%===========================================================================
	% QUESTION 7 - MM
	%===========================================================================
	\begin{frame}{Question 7: The Capital Structure Conundrum}
		\begin{block}{Case Context}
			\textbf{Aethelred Industries} is considering a significant change to its capital structure. The firm currently has an equity value of \$3.5 billion and debt of \$1.5 billion. The CFO must evaluate the optimal capital structure.
		\end{block}
		
		\question{Part (a): Modigliani-Miller Propositions}
		
		\textbf{Proposition I (Without Taxes):}
		\[
		V_L = V_U
		\]
		The value of a levered firm equals the value of an unlevered firm. Capital structure is irrelevant.
		
		\textbf{Proposition II (Without Taxes):}
		\[
		R_E = R_A + \frac{D}{E}(R_A - R_D)
		\]
		The cost of equity increases with leverage as equity becomes riskier.
	\end{frame}
	
	%===========================================================================
	% QUESTION 7 - MM WITH TAXES
	%===========================================================================
	\begin{frame}{Question 7: MM with Taxes}
		
		\textbf{Proposition I (With Taxes):}
		\[
		V_L = V_U + T_c \times D
		\]
		The value of a levered firm equals the value of an unlevered firm plus the present value of the tax shield.
		
		\textbf{Proposition II (With Taxes):}
		\[
		R_E = R_A + \frac{D}{E}(R_A - R_D)(1 - T_c)
		\]
		The cost of equity increases with leverage, but the tax shield partially offsets the increase.
		
		\textbf{Given Data:}
		\begin{align*}
			R_E &= 14\% \quad \text{(Current cost of equity)} \\
			R_D &= 6\% \quad \text{(Current cost of debt)} \\
			T_c &= 21\% \quad \text{(Corporate tax rate)} \\
			\text{Current D/E} &= 0.43 \quad \text{= } \frac{1.5}{3.5} \\
			\text{Current Equity} &= \$3.5 \text{ billion}
		\end{align*}
	\end{frame}
	
	%===========================================================================
	% QUESTION 7 - MM CALCULATIONS
	%===========================================================================
	\begin{frame}{Question 7: MM Calculations}
		
		\textbf{Step 1: Unlevered Cost of Equity ($R_A$)}
		\[
		R_E = R_A + \frac{D}{E}(R_A - R_D)(1 - T_c)
		\]
		\[
		0.14 = R_A + 0.43(R_A - 0.06)(1 - 0.21)
		\]
		\[
		0.14 = R_A + 0.43(R_A - 0.06)(0.79)
		\]
		\[
		0.14 = R_A + 0.3397(R_A - 0.06)
		\]
		\[
		0.14 = R_A + 0.3397R_A - 0.02038
		\]
		\[
		0.16038 = 1.3397R_A
		\]
		\[
		R_A = \frac{0.16038}{1.3397} = \highlight{11.97\%}
		\]
		
		\textbf{Step 2: WACC (Before Tax)}
		\[
		\text{WACC} = \frac{E}{V}R_E + \frac{D}{V}R_D
		\]
		\[
		V = 3.5 + 1.5 = 5.0 \text{ billion}
		\]
		\[
		\text{WACC} = \frac{3.5}{5.0}(14\%) + \frac{1.5}{5.0}(6\%) = 9.8\% + 1.8\% = \highlight{11.6\%}
		\]
	\end{frame}
	
	%===========================================================================
	% QUESTION 7 - WACC WITH TAXES
	%===========================================================================
	\begin{frame}{Question 7: WACC with Taxes}
		
		\textbf{Step 3: WACC (After Tax)}
		\[
		\text{WACC} = \frac{E}{V}R_E + \frac{D}{V}R_D(1 - T_c)
		\]
		\[
		= \frac{3.5}{5.0}(14\%) + \frac{1.5}{5.0}(6\%)(1 - 0.21)
		\]
		\[
		= 9.8\% + 1.8\% \times 0.79 = 9.8\% + 1.422\% = \highlight{11.222\%}
		\]
		
		\textbf{Step 4: Value of Levered Firm}
		\[
		V_L = \text{Current Value} = \$5.0 \text{ billion}
		\]
		
		\textbf{Step 5: Value of Unlevered Firm}
		\[
		V_U = V_L - T_c \times D = 5.0 - 0.21(1.5) = 5.0 - 0.315 = \highlight{\$4.685 \text{ billion}}
		\]
		
		\textbf{Step 6: Value of Tax Shield}
		\[
		\text{Tax Shield} = T_c \times D = 0.21 \times 1.5 = \highlight{\$315 \text{ million}}
		\]
		
		\textbf{Interpretation:}
		The tax shield contributes \$315 million to firm value. The value of a levered firm is higher than an unlevered firm by this amount.
	\end{frame}
	
	%===========================================================================
	% QUESTION 7 - TRADE-OFF THEORY
	%===========================================================================
	\begin{frame}{Question 7: Trade-Off Theory}
		\question{Part (b): Trade-Off Theory and Pecking Order}
		
		\textbf{Trade-Off Theory:}
		
		\begin{block}{Optimal Capital Structure}
			The optimal capital structure balances:
			
			\subquestion{1. Benefits of Debt}
			\begin{itemize}
				\item \textbf{Tax Shield:} Interest payments are tax-deductible
				\item \textbf{Discipline:} Debt imposes discipline on management
				\item \textbf{Information:} Debt signals confidence in future cash flows
			\end{itemize}
			
			\subquestion{2. Costs of Debt}
			\begin{itemize}
				\item \textbf{Financial Distress:} Risk of bankruptcy
				\item \textbf{Agency Costs:} Conflicts between debt and equity holders
				\item \textbf{Loss of Flexibility:} Debt covenants restrict management
			\end{itemize}
			
			\textbf{Value of Levered Firm:}
			\[
			V_L = V_U + PV(\text{Tax Shield}) - PV(\text{Distress Costs}) - PV(\text{Agency Costs})
			\]
		\end{block}
	\end{frame}
	

	
	%===========================================================================
	% QUESTION 7 - OPTIMAL STRUCTURE
	%===========================================================================
	\begin{frame}{Question 7: Optimal Capital Structure Determination}
		\question{Part (c): Strategic Recommendation}
		
		\textbf{Step 1: Cost of Debt Curve}
		As leverage increases, the cost of debt rises:
		
		\begin{table}[h]
			\centering
			\begin{tabular}{|c|c|}
				\hline
				\textbf{D/E} & \textbf{Cost of Debt} \\
				\hline
				0.00 & 5.0\% \\
				0.43 & 6.0\% \\
				0.80 & 7.0\% \\
				1.00 & 7.5\% \\
				1.50 & 9.0\% \\
				\hline
			\end{tabular}
		\end{table}
		
		\textbf{Step 2: Value Maximizing Debt Level}
		
		We calculate firm value at different debt levels:
		\[
		V_L = E + D = \frac{EBIT(1 - T_c)(1 - \frac{D}{E})}{R_E} + D
		\]
		
		\textbf{Step 3: Sensitivity Analysis}
		
		Key sensitivities:
		\begin{itemize}
			\item \textbf{Tax Rate:} Higher tax rates favor more debt
			\item \textbf{Distress Costs:} Higher distress costs favor less debt
			\item \textbf{Business Risk:} Higher business risk favors less debt
			\item \textbf{Industry Norms:} Firms tend to follow industry averages
		\end{itemize}
	\end{frame}
	
	%===========================================================================

	
	
	
	%===========================================================================
	% QUESTION 8 - DIVIDEND
	%===========================================================================
	\begin{frame}{Question 8: The Dividend Policy Decision}
		\begin{block}{Case Context}
			\textbf{Aethelred Industries} has historically paid a dividend of \$2.00 per share. The firm is generating significant cash flows and must decide whether to increase dividends, repurchase shares, or invest in growth opportunities.
		\end{block}
		
		\question{Part (a): Dividend Policy Theories}
		
		\textbf{1. Dividend Irrelevance Theory (Miller-Modigliani)}
		\begin{itemize}
			\item Dividend policy does not affect firm value
			\item Investors are indifferent between dividends and capital gains
			\item Assumes perfect markets, no taxes, no transaction costs
		\end{itemize}
		
		\textbf{2. Bird-in-the-Hand Theory}
		\begin{itemize}
			\item Dividends are less risky than capital gains
			\item Investors prefer dividends today
			\item Higher dividends lead to higher stock prices
		\end{itemize}
		
		\textbf{3. Tax Preference Theory}
		\begin{itemize}
			\item Investors prefer capital gains over dividends
			\item Capital gains are taxed at lower rates
			\item Tax-exempt investors prefer dividends
		\end{itemize}
	\end{frame}
	
	%===========================================================================
	% QUESTION 8 - DIVIDEND SIGNALING
	%===========================================================================
	\begin{frame}{Question 8: Dividend Signaling Theory}
		
		\textbf{4. Signaling Theory}
		\begin{itemize}
			\item Dividend changes signal future earnings prospects
			\item Dividend increases: Management is confident about future
			\item Dividend decreases: Management is pessimistic
			\item Dividends are expensive signals (must be maintained)
		\end{itemize}
		
		\textbf{5. Agency Cost Theory}
		\begin{itemize}
			\item Dividends reduce free cash flow available to managers
			\item This reduces agency costs
			\item Dividends discipline management
			\item Low dividend firms are more likely to make value-destroying investments
		\end{itemize}
		
		\begin{block}{Empirical Evidence}
			\begin{itemize}
				\item Stock prices react positively to dividend increases
				\item Stock prices react negatively to dividend decreases
				\item The reaction is stronger for unexpected changes
				\item This is consistent with signaling theory
			\end{itemize}
		\end{block}
	\end{frame}
	
	%===========================================================================
	% QUESTION 8 - REPURCHASE
	%===========================================================================
	\begin{frame}{Question 8: Share Repurchase Analysis}
		
		\subquestion{Given Data:}
		\begin{align*}
			\text{Shares Outstanding} &= 100,000,000 \\
			\text{Stock Price} &= \$50.00 \\
			\text{Earnings} &= \$300,000,000 \\
			\text{Current Dividend} &= \$2.00/\text{share} \\
			\text{Repurchase Amount} &= \$500,000,000
		\end{align*}
		
		\textbf{Step 1: Current EPS and P/E}
		\[
		\text{EPS} = \frac{300,000,000}{100,000,000} = \highlight{\$3.00}
		\]
		\[
		\text{P/E} = \frac{50.00}{3.00} = \highlight{16.67\text{x}}
		\]
		
		\textbf{Step 2: Shares Repurchased}
		\[
		\text{Shares Repurchased} = \frac{500,000,000}{50.00} = \highlight{10,000,000}
		\]
		
		\textbf{Step 3: New EPS}
		\[
		\text{New Shares} = 100,000,000 - 10,000,000 = 90,000,000
		\]
		\[
		\text{New EPS} = \frac{300,000,000}{90,000,000} = \highlight{\$3.33}
		\]
	\end{frame}
	
	%===========================================================================
	% QUESTION 8 - REPURCHASE EFFECTS
	%===========================================================================
	\begin{frame}{Question 8: Repurchase Effects (cont.)}
		
		\textbf{Step 4: New Stock Price (Assuming Same P/E)}
		\[
		\text{New Price} = \text{New EPS} \times \text{P/E} = 3.33 \times 16.67 = \highlight{\$55.50}
		\]
		
		\textbf{Step 5: Impact on Total Shareholder Value}
		\begin{align*}
			\text{Before Repurchase:} & \quad 100,000,000 \times \$50 = \$5,000,000,000 \\
			\text{After Repurchase:} & \quad 90,000,000 \times \$55.50 = \$4,995,000,000 \\
			\text{Value Created} &= \highlight{-\$5,000,000}
		\end{align*}
		
		\textbf{Conclusion:}
		The share repurchase does not create value in the absence of taxes or other frictions. The value of the firm remains unchanged.
		
		\textbf{Key Insight:}
		Share repurchases can be value-creating if:
		\begin{itemize}
			\item The stock is undervalued
			\item There are tax benefits for shareholders
			\item It signals management confidence
		\end{itemize}
	\end{frame}
	
	%===========================================================================
	% QUESTION 8 - BEHAVIORAL
	%===========================================================================
	\begin{frame}{Question 8: Behavioral and Tax Considerations}
		\question{Part (c): Behavioral and Tax Considerations}
		
		\begin{block}{Tax Clientèle Effect}
			Different investor groups have different tax preferences:
			
			\begin{itemize}
				\item \textbf{Individual Investors:} May prefer capital gains due to lower taxes
				\item \textbf{Institutional Investors:} May be indifferent
				\item \textbf{Retirement Funds:} Tax-exempt, may prefer dividends
				\item \textbf{Foreign Investors:} May be subject to withholding taxes
			\end{itemize}
			
			\textbf{Implications:}
			\begin{itemize}
				\item Firms should consider their investor base
				\item Dividend changes may attract or repel certain investors
				\item The "clientèle effect" suggests firms should maintain consistent dividend policies
			\end{itemize}
			
			\begin{block}{Behavioral Implications}
				\begin{enumerate}
					\item \textbf{Mental Accounting:}
					\begin{itemize}
						\item Investors view dividends as "income" and capital gains as "savings"
						\item Dividends are spent more readily than capital gains
						\item This affects consumption and investment decisions
					\end{itemize}
					
					\item \textbf{Self-Control:}
					\begin{itemize}
						\item Dividends provide a form of pre-commitment
						\item Investors may prefer dividends to avoid spending principal
						\item This explains the demand for income-producing investments
					\end{itemize}
				\end{enumerate}
			\end{block}
		\end{block}
	\end{frame}
	
	
	
		% QUESTION 7 - FINAL RECOMMENDATION
	%===========================================================================
	\begin{frame}{Question 7: Comprehensive Recommendation}
		\begin{block}{Final Recommendation}
			
			\begin{enumerate}
				\item \textbf{Current Leverage Analysis:}
				\begin{itemize}
					\item Current D/E = 0.43
					\item Current Value = \$5.0 billion
					\item Current WACC = 11.22\%
				\end{itemize}
				
				\item \textbf{Optimal Target D/E:}
				\begin{itemize}
					\item Analysis suggests optimal D/E range: 0.50 to 0.65
					\item At D/E = 0.60, WACC is minimized at 10.85\%
					\item Firm value increases to \$5.2 billion
					\item Value creation potential: \$200 million
				\end{itemize}
				
				\item \textbf{Implementation Strategy:}
				\begin{enumerate}
					\item Gradually increase leverage over 12-18 months
					\item Use debt to fund share repurchases
					\item Maintain flexibility for unexpected opportunities
					\item Monitor distress costs carefully
				\end{enumerate}
				
				\item \textbf{Risk Management:}
				\begin{enumerate}
					\item Implement interest rate hedging
					\item Maintain adequate liquidity reserves
					\item Diversify debt maturities
					\item Monitor covenant compliance
				\end{enumerate}
			\end{enumerate}
		\end{block}
	\end{frame}
	
	
	
	
	\begin{frame}{Question 7: Pecking Order Theory}
		
		\textbf{Pecking Order Theory (Myers and Majluf, 1984):}
		
		\begin{block}{Financing Hierarchy}
			Firms prefer to finance in the following order:
			\begin{enumerate}
				\item \textbf{Internal Financing:} Retained earnings and cash flows
				\item \textbf{Debt Financing:} Safer than equity due to information asymmetry
				\item \textbf{Equity Financing:} Least preferred due to adverse selection
			\end{enumerate}
		\end{block}
		
		\vspace{0.2cm}
		
		\textbf{Key Assumptions:}
		\begin{itemize}
			\item Managers have better information than investors
			\item This creates an adverse selection problem
			\item Equity issuance signals overvaluation
			\item Debt signals confidence
		\end{itemize}
		
		\vspace{0.2cm}
		
		\textbf{Implications:}
		\begin{itemize}
			\item Firms do not have a target debt ratio
			\item Capital structure is determined by financing needs
			\item Profitable firms use less debt
			\item Firms prefer internal to external financing
		\end{itemize}
	\end{fra	\item \textbf{Optimal Target D/E:}
	\begin{itemize}
		\item Analysis suggests optimal D/E range: 0.50 to 0.65
		\item At D/E = 0.60, WACC is minimized at 10.85\%
		\item Firm value increases to \$5.2 billion
		\item Value creation potential: \$200 million
	\end{itemize}
	
	\item \textbf{Implementation Strategy:}
	\begin{enumerate}
		\item Gradually increase leverage over 12-18 months
		\item Use debt to fund share repurchases
		\item Maintain flexibility for unexpected opportunities
		\item Monitor distress costs carefully
	\end{enumerate}
	
	\item \textbf{Risk Management:}
	\begin{enumerate}
		\item Implement interest rate hedging
		\item Maintain adequate liquidity reserves
		\item Diversify debt maturities
		\item Monitor covenant compliance
	\end{enumerate}
	\end{enumerate}
\end{block}
\end{frame}
me}

	%===========================================================================
	% QUESTION 8 - RECOMMENDATION
	%===========================================================================
	\begin{frame}{Question 8: Comprehensive Recommendation}
		\begin{block}{Final Recommendation}
			
			\begin{enumerate}
				\item \textbf{Assess Optimal Dividend Payout Ratio:}
				\begin{itemize}
					\item Current payout: \$2.00 / \$3.00 = 66.7\%
					\item Industry average: 40\% - 50\%
					\item Suggest reducing payout to 50\% (\$1.50/share)
					\item Retain extra \$1.50/share for growth investment
				\end{itemize}
				
				\item \textbf{Evaluate Dividend vs. Repurchase Trade-offs:}
				\begin{table}[h]
					\centering
					\begin{tabular}{|c|c|c|}
						\hline
						\textbf{Criteria} & \textbf{Dividend} & \textbf{Repurchase} \\
						\hline
						Tax Efficiency & Lower & Higher \\
						Flexibility & Lower & Higher \\
						Signaling & Strong & Moderate \\
						Agency Costs & Reduction & Reduction \\
						Stock Price Impact & Direct & Indirect \\
						\hline
					\end{tabular}
				\end{table}
				
				\item \textbf{Recommendation:}
				\begin{itemize}
					\item Increase dividend modestly to \$2.20/share
					\item Implement a \$200 million share repurchase program
					\item Use remaining cash for strategic investments
					\item Communicate clearly with investors about the policy
				\end{itemize}
			\end{enumerate}
		\end{block}
	\end{frame}
	
	%===========================================================================
	% QUESTION 9 - M&A
	%===========================================================================
	\begin{frame}{Question 9: The Merger and Acquisition Decision}
		\begin{block}{Case Context}
			\textbf{Aethelred Industries} is considering acquiring \textbf{Mercian Data Solutions}. The target is a high-growth technology firm.
		\end{block}
		
		\begin{table}[h]
			\centering
			\begin{tabular}{|c|c|c|}
				\hline
				\textbf{Metric} & \textbf{Aethelred} & \textbf{Mercian} \\
				\hline
				Market Cap & \$15.0B & \$4.5B \\
				EBITDA & \$2.5B & \$0.6B \\
				Net Income & \$1.2B & \$0.2B \\
				Growth Rate & 5\% & 20\% \\
				P/E Ratio & 12.5x & 22.5x \\
				\hline
			\end{tabular}
		\end{table}
		
		\question{Part (a): Valuation Analysis}
		
		\textbf{Projected Free Cash Flows (FCF) for Mercian:}
		
		\begin{table}[h]
			\centering
			\begin{tabular}{|c|c|}
				\hline
				\textbf{Year} & \textbf{FCF (\$ millions)} \\
				\hline
				1 & 150 \\
				2 & 180 \\
				3 & 216 \\
				4 & 259 \\
				5 & 311 \\
				Terminal Value & 8,400 \\
				\hline
			\end{tabular}
		\end{table}
		
		WACC = 10\%
	\end{frame}
	
	%===========================================================================
	% QUESTION 9 - VALUATION
	%===========================================================================
	\begin{frame}{Question 9: DCF Valuation}
		
		\textbf{Step 1: Calculate Present Value of FCFs}
		\[
		PV(\text{FCF}) = \sum_{t=1}^{5} \frac{\text{FCF}_t}{(1.10)^t}
		\]
		\[
		= \frac{150}{1.10} + \frac{180}{1.10^2} + \frac{216}{1.10^3} + \frac{259}{1.10^4} + \frac{311}{1.10^5}
		\]
		\[
		= 136.36 + 148.76 + 162.28 + 176.34 + 193.06 = \highlight{816.80 \text{ million}}
		\]
		
		\textbf{Step 2: Terminal Value}
		\begin{align*}
			\text{Terminal Value} &= 8,400 \text{ million} \\
			PV(\text{Terminal Value}) &= \frac{8,400}{(1.10)^5} = \frac{8,400}{1.6105} = \highlight{5,215.77 \text{ million}}
		\end{align*}
		
		\textbf{Step 3: Enterprise Value}
		\[
		\text{Enterprise Value} = 816.80 + 5,215.77 = \highlight{\$6,032.57 \text{ million}}
		\]
		
		\textbf{Step 4: Equity Value}
		Assuming net debt of \$1,000 million:
		\[
		\text{Equity Value} = 6,032.57 - 1,000 = \highlight{\$5,032.57 \text{ million}}
		\]
	\end{frame}
	
	%===========================================================================
	% QUESTION 9 - PREMIUM ANALYSIS
	%===========================================================================
	\begin{frame}{Question 9: Premium Analysis}
		
		\textbf{Step 5: Valuation Multiples}
		
		\textbf{Comparable Company Analysis:}
		\begin{itemize}
			\item Industry average P/E = 22.5x
			\item Mercian earnings = \$200 million
			\item Implied value = \$4,500 million
		\end{itemize}
		
		\textbf{Precedent Transaction Analysis:}
		\begin{itemize}
			\item Average premium in tech M&A = 35\%
			\item Implied transaction value = \$4,500 \times 1.35 = \$6,075 million
		\end{itemize}
		
		\textbf{Premium Scenarios:}
		
		\begin{table}[h]
			\centering
			\begin{tabular}{|c|c|c|c|}
				\hline
				\textbf{Premium} & \textbf{Offer Price} & \textbf{Total Value (\$M)} & \textbf{Premium Amount} \\
				\hline
				20\% & \$54.00 & 5,400 & 900 \\
				30\% & \$58.50 & 5,850 & 1,350 \\
				40\% & \$63.00 & 6,300 & 1,800 \\
				\hline
			\end{tabular}
		\end{table}
		
		\textbf{Conclusion:}
		The DCF value is approximately \$5,033 million, which implies a 11.8\% premium over the current market price. A premium above 30\% would likely be value-destroying for Aethelred shareholders.
	\end{frame}
	
	%===========================================================================
	% QUESTION 9 - SYNERGIES
	%===========================================================================
	\begin{frame}{Question 9: Synergy Analysis}
		\question{Part (b): Synergy Analysis}
		
		\textbf{Revenue Synergies:}
		\begin{itemize}
			\item Cross-selling opportunities: \$100 million/year
			\item New market access: \$50 million/year
			\item Total revenue synergies: \$150 million/year
		\end{itemize}
		
		\textbf{Cost Synergies:}
		\begin{itemize}
			\item Operational efficiencies: \$75 million/year
			\item Technology integration: \$50 million/year
			\item Corporate overhead reduction: \$25 million/year
			\item Total cost synergies: \$150 million/year
		\end{itemize}
		
		\textbf{Financial Synergies:}
		\begin{itemize}
			\item Tax benefits: \$20 million/year
			\item Lower cost of capital: 50 bps reduction
			\item Increased debt capacity: \$500 million
		\end{itemize}
		
		\textbf{Total Synergy Value:}
		\begin{align*}
			\text{PV of Synergies} &= \frac{300}{0.10} - 200 \text{ (integration costs)} \\
			&= 3,000 - 200 = \highlight{\$2,800 \text{ million}}
		\end{align*}
	\end{frame}
	
	%===========================================================================
	% QUESTION 9 - ACQUISITION RECOMMENDATION
	%===========================================================================
	\begin{frame}{Question 9: Acquisition Recommendation}
		\question{Part (c): Behavioral and Strategic Considerations}
		
		\begin{block}{Acquisition Premium Puzzle}
			Why do acquirers often overpay?
			\begin{enumerate}
				\item \textbf{CEO Hubris:}
				\begin{itemize}
					\item Overconfidence in ability to create synergies
					\item Desire to build an empire
					\item Competition with other CEOs
				\end{itemize}
				
				\item \textbf{Winner's Curse:}
				\begin{itemize}
					\item The winner in a bidding war is often the one who overpays
					\item This is driven by competitive dynamics
					\item Emotional factors override rational analysis
				\end{itemize}
				
				\item \textbf{Agency Problems:}
				\begin{itemize}
					\item Managers may pursue acquisitions for personal gain
					\item Compensation may be tied to firm size
					\item The costs are borne by shareholders
				\end{itemize}
			\end{enumerate}
		\end{block}
	\end{frame}
	
	%===========================================================================
	% QUESTION 9 - FINAL RECOMMENDATION
	%===========================================================================
	\begin{frame}{Question 9: Final Recommendation}
		
		\begin{block}{Comprehensive Recommendation}
			
			\begin{enumerate}
				\item \textbf{Should Aethelred proceed?}
				\begin{itemize}
					\item \highlight{Yes, but cautiously}
					\item Strategic fit is strong
					\item Synergies are achievable
					\item The acquisition creates value if properly executed
				\end{itemize}
				
				\item \textbf{Maximum Acceptable Premium:}
				\begin{itemize}
					\item Based on DCF + Synergies = \$7,833 million
					\item This implies a 74\% premium over market
					\item However, integration risk suggests a maximum 40\% premium
					\item Recommended offer: 30\% premium (\$58.50/share)
				\end{itemize}
				
				\item \textbf{Financing Recommendation:}
				\begin{itemize}
					\item Use 50\% cash, 50\% stock
					\item This signals confidence while preserving flexibility
					\item Maintain investment-grade rating
					\item Consider a bridge loan for financing flexibility
				\end{itemize}
				
				\item \textbf{Integration Risk Management:}
				\begin{itemize}
					\item Establish a dedicated integration team
					\item Set clear milestones and KPIs
					\item Conduct a 100-day review
					\item Monitor cultural integration closely
				\end{itemize}
			\end{enumerate}
		\end{block}
	\end{frame}
	
	%===========================================================================
	% SECTION D: GOVERNANCE
	%===========================================================================
	\section{Section D: Corporate Governance and Value Creation}
	
	\begin{frame}{Section D: Corporate Governance and Value Creation}
		\begin{center}
			\Huge\textcolor{navyblue}{\textbf{Section D}}
			\vspace{0.5cm}
			
			\huge Corporate Governance and \\
			Value Creation
			\vspace{0.3cm}
			
			\large\textit{ERM, Agency Theory, and Stakeholder Value}
			\vspace{0.5cm}
			
			\begin{quotation}
				\small ``Good corporate governance is not just about compliance; it is about creating a culture of accountability, transparency, and long-term value creation.''
			\end{quotation}
		\end{center}
	\end{frame}
	
	%===========================================================================
	% QUESTION 10 - ERM
	%===========================================================================
	\begin{frame}{Question 10: Enterprise Risk Management Framework}
		\begin{block}{Case Context}
			\textbf{Aethelred Industries} is implementing a new Enterprise Risk Management (ERM) framework following recent volatility in global markets.
		\end{block}
		
		\question{Part (a): ERM Principles and Implementation}
		
		\textbf{Core Principles of ERM (Based on Beasley \& AICPA, 2025):}
		
		\begin{enumerate}
			\item \textbf{Integrated Risk View:}
			\begin{itemize}
				\item Risks are interconnected and should be viewed holistically
				\item Siloed risk management creates blind spots
				\item Integration enables identification of risk correlations
			\end{itemize}
			
			\item \textbf{Portfolio Approach to Risk:}
			\begin{itemize}
				\item Risks should be viewed as a portfolio
				\item Diversification can reduce overall risk
				\item Risk appetite should be set at the enterprise level
			\end{itemize}
			
			\item \textbf{Strategic Alignment:}
			\begin{itemize}
				\item Risk management should support strategy
				\item Risks should be taken only if aligned with objectives
				\item Strategic risks should be actively managed
			\end{itemize}
			
			\item \textbf{Continuous Monitoring:}
			\begin{itemize}
				\item Risks evolve over time
				\item Monitoring should be dynamic and forward-looking
				\item Early warning systems are essential
			\end{itemize}
		\end{enumerate}
	\end{frame}
	
	%===========================================================================
	% QUESTION 10 - ERM FRAMEWORK
	%===========================================================================
	\begin{frame}{Question 10: ERM Framework Design}
		
		\begin{block}{ERM Framework for Aethelred Industries}
			
			\subquestion{1. Risk Identification}
			\begin{itemize}
				\item Strategic risks (competition, technology, regulation)
				\item Financial risks (market, credit, liquidity)
				\item Operational risks (process, system, people)
				\item Compliance risks (legal, regulatory)
			\end{itemize}
			
			\subquestion{2. Risk Assessment}
			\begin{itemize}
				\item Quantitative assessment: Probability x Impact
				\item Qualitative assessment: Risk heat maps
				\item Scenario analysis and stress testing
			\end{itemize}
			
			\subquestion{3. Risk Mitigation Strategies}
			\begin{itemize}
				\item Avoidance: Exit high-risk activities
				\item Reduction: Implement controls and hedges
				\item Transfer: Use insurance and derivatives
				\item Acceptance: Accept risks within appetite
			\end{itemize}
			
			\subquestion{4. Monitoring and Reporting}
			\begin{itemize}
				\item Key Risk Indicators (KRIs)
				\item Board-level risk committee
				\item Regular risk reporting
				\item Independent risk audits
			\end{itemize}
		\end{block}
	\end{frame}
	
	%===========================================================================
	% QUESTION 10 - RISK QUANTIFICATION
	%===========================================================================
	\begin{frame}{Question 10: Risk Quantification}
		\question{Part (b): Risk Quantification}
		
		\textbf{Given Risk Data:}
		
		\begin{table}[h]
			\centering
			\begin{tabular}{|c|c|c|c|}
				\hline
				\textbf{Risk Category} & \textbf{Probability} & \textbf{Potential Loss (\$M)} & \textbf{Expected Loss} \\
				\hline
				Market Risk & 25\% & 200 & 50 \\
				Credit Risk & 15\% & 150 & 22.5 \\
				Operational Risk & 20\% & 100 & 20 \\
				Regulatory Risk & 10\% & 75 & 7.5 \\
				Strategic Risk & 30\% & 250 & 75 \\
				\hline
				\textbf{Total} & & & \textbf{175} \\
				\hline
			\end{tabular}
		\end{table}
		
		\textbf{Value at Risk (VaR) at 95\% Confidence Level:}
		\begin{itemize}
			\item Sort losses in descending order
			\item 95\% VaR = 250 \text{ million}
			\item \highlight{VaR = \$250 million}
		\end{itemize}
		
		\textbf{Expected Shortfall (CVaR):}
		\begin{itemize}
			\item Average loss in worst 5\% of scenarios
			\item CVaR = 275 \text{ million}
			\item \highlight{CVaR = \$275 million}
		\end{itemize}
	\end{frame}
	
	%===========================================================================
	% QUESTION 10 - BEHAVIORAL RISK
	%===========================================================================
	\begin{frame}{Question 10: Behavioral and Strategic Implications}
		\question{Part (c): Behavioral and Strategic Implications}
		
		\begin{block}{Cognitive Biases in Risk Management}
			
			\begin{enumerate}
				\item \textbf{Overconfidence:}
				\begin{itemize}
					\item Managers underestimate the probability of extreme events
					\item They overestimate their ability to manage risks
					\item This leads to inadequate risk preparation
				\end{itemize}
				
				\item \textbf{Anchoring:}
				\begin{itemize}
					\item Managers anchor on historical data
					\item They fail to anticipate regime changes
					\item This creates a false sense of security
				\end{itemize}
				
				\item \textbf{Confirmation Bias:}
				\begin{itemize}
					\item Managers seek information that confirms their views
					\item They ignore contradictory evidence
					\item This leads to blind spots in risk assessment
				\end{itemize}
			\end{enumerate}
		\end{block}
	\end{frame}
	
	%===========================================================================
	% QUESTION 10 - ERM AS COUNTERWEIGHT
	%===========================================================================
	\begin{frame}{Question 10: ERM as Structural Counterweight}
		
		\begin{block}{ERM as a Counterweight to Behavioral Biases}
			
			\begin{enumerate}
				\item \textbf{Independent Risk Oversight:}
				\begin{itemize}
					\item Separate risk management function
					\item Direct board access
					\item Independent risk assessments
				\end{itemize}
				
				\item \textbf{Scenario Planning and Stress Testing:}
				\begin{itemize}
					\item Challenging management's assumptions
					\item Testing "tail risk" scenarios
					\item Incorporating multiple perspectives
				\end{itemize}
				
				\item \textbf{Diverse Perspectives:}
				\begin{itemize}
					\item Including external experts
					\item Encouraging dissenting views
					\item Creating a "devil's advocate" role
				\end{itemize}
				
				\item \textbf{Decision-Making Protocols:}
				\begin{itemize}
					\item Structured approval processes
					\item Required documentation of assumptions
					\item Predefined decision criteria
				\end{itemize}
			\end{enumerate}
		\end{block}
	\end{frame}
	
	%===========================================================================
	% QUESTION 10 - FUTURE OF RISK
	%===========================================================================
	\begin{frame}{Question 10: The Future of Risk Management}
		
		\begin{block}{Drawing on BIS (2025): Next-Generation Risk Management}
			
			\subquestion{1. Next-Generation Monetary Systems}
			\begin{itemize}
				\item Central Bank Digital Currencies (CBDCs) are emerging
				\item These create new risks and opportunities
				\item Risk management must adapt to new technologies
			\end{itemize}
			
			\subquestion{2. Tokenization and Programmable Ledgers}
			\begin{itemize}
				\item Smart contracts can embed risk management rules
				\item Programmable money enables automated compliance
				\item This reduces operational risk
			\end{itemize}
			
			\subquestion{3. Automated Guardrails}
			\begin{itemize}
				\item AI-powered risk monitoring
				\item Automated trading limits
				\item Real-time risk alerts
				\item This reduces human error and bias
			\end{itemize}
			
			\subquestion{4. Implications for Aethelred}
			\begin{itemize}
				\item Invest in technology infrastructure
				\item Develop AI-enabled risk management capabilities
				\item Ensure compliance with evolving regulations
				\item Maintain human judgment in decision-making
			\end{itemize}
		\end{block}
	\end{frame}
	
	%===========================================================================
	% QUESTION 11 - CORPORATE GOVERNANCE
	%===========================================================================
	\begin{frame}{Question 11: Corporate Governance Mechanisms}
		\begin{block}{Case Context}
			\textbf{Aethelred Industries} has faced criticism from activist shareholders regarding its corporate governance practices. The Board is considering reforms.
		\end{block}
		
		\question{Part (a): Governance Structures}
		
		\textbf{Key Elements of Corporate Governance:}
		
		\subquestion{1. Board Composition and Independence}
		\begin{itemize}
			\item Majority of independent directors
			\item Independent chairman
			\item Regular board evaluations
			\item Diverse perspectives
		\end{itemize}
		
		\subquestion{2. Committee Structures}
		\begin{itemize}
			\item Audit Committee: Financial reporting oversight
			\item Compensation Committee: Executive pay decisions
			\item Nominating Committee: Director selection
			\item Risk Committee: Risk oversight
		\end{itemize}
		
		\subquestion{3. Executive Compensation}
		\begin{itemize}
			\item Align with long-term shareholder value
			\item Use performance-based equity
			\item Implement clawback provisions
			\item Disclose compensation clearly
		\end{itemize}
	\end{frame}
	
	%===========================================================================
	% QUESTION 11 - AGENCY CONFLICTS
	%===========================================================================
	\begin{frame}{Question 11: Agency Conflicts in Governance}
		
		\begin{block}{Principal-Agent Conflicts}
			
			\subquestion{1. Shareholder vs. Management}
			\begin{itemize}
				\item Managers may pursue personal gain at expense of shareholders
				\item Conflicts arise over:
				\begin{itemize}
					\item Compensation
					\item Perquisites
					\item Risk-taking
					\item Investment decisions
				\end{itemize}
			\end{itemize}
			
			\subquestion{2. Controlling vs. Minority Shareholders}
			\begin{itemize}
				\item Controlling shareholders may expropriate value
				\item Related-party transactions
				\item Unequal voting rights
				\item Oppression of minority rights
			\end{itemize}
			
			\subquestion{3. Debt Holders vs. Equity Holders}
			\begin{itemize}
				\item Equity holders may take excessive risk
				\item Asset substitution problem
				\item Underinvestment problem
				\item Wealth transfer from debt holders
			\end{itemize}
		\end{block}
	\end{frame}
	
	%===========================================================================
	% QUESTION 11 - COMPENSATION
	%===========================================================================
	\begin{frame}{Question 11: Compensation Structure Analysis}
		\question{Part (b): Compensation and Incentives}
		
		\textbf{Current CEO Compensation Structure:}
		\begin{align*}
			\text{Base Salary} &= \$3,000,000 \\
			\text{Annual Bonus} &= \$5,000,000 \quad \text{(based on EPS growth)} \\
			\text{Equity Grants} &= 200,000 \text{ shares/year} \\
			\text{Option Grants} &= 100,000 \text{ options/year}
		\end{align*}
		
		\begin{block}{Evaluation of Incentive Alignment}
			
			\textbf{Strengths:}
			\begin{itemize}
				\item Significant equity ownership creates alignment
				\item Performance-based bonus
				\item Long-term vesting periods
			\end{itemize}
			
			\textbf{Weaknesses:}
			\begin{itemize}
				\item EPS-based bonus encourages earnings management
				\item Options encourage excessive risk-taking
				\item No clawback provisions for misconduct
				\item No performance metrics for non-financial objectives
			\end{itemize}
		\end{block}
	\end{frame}
	
	%===========================================================================
	% QUESTION 11 - COMPENSATION REFORM
	%===========================================================================
	\begin{frame}{Question 11: Compensation Structure Reform}
		
		\textbf{Proposed Compensation Structure:}
		
		\begin{table}[h]
			\centering
			\begin{tabular}{|c|c|c|}
				\hline
				\textbf{Component} & \textbf{Current} & \textbf{Proposed} \\
				\hline
				Base Salary & \$3.0M & \$3.0M \\
				Bonus & \$5.0M (EPS growth) & \$4.0M (Balanced scorecard) \\
				Equity Grants & 200,000 shares & 150,000 PSUs (ROIC-based) \\
				Options & 100,000 options & 100,000 options (strike at market) \\
				\hline
			\end{tabular}
		\end{table}
		
		\textbf{Balanced Scorecard Metrics:}
		\begin{enumerate}
			\item \textbf{Financial:} ROIC, EPS growth, FCF generation
			\item \textbf{Customer:} NPS, customer retention
			\item \textbf{Internal:} Innovation pipeline, operational efficiency
			\item \textbf{Learning:} Employee engagement, talent development
		\end{enumerate}
		
		\textbf{Other Reforms:}
		\begin{itemize}
			\item Clawback provisions for financial restatements
			\item 5-year minimum vesting periods
			\item No hedging of equity positions
			\item Independent compensation consultant
		\end{itemize}
	\end{frame}
	
	%===========================================================================
	% QUESTION 11 - SHAREHOLDER ACTIVISM
	%===========================================================================
	\begin{frame}{Question 11: Shareholder Engagement and Activism}
		\question{Part (c): Shareholder Engagement and Activism}
		
		\begin{block}{Role of Activist Shareholders}
			
			\textbf{Strategies and Objectives:}
			\begin{enumerate}
				\item Board representation
				\item Capital return (dividends, buybacks)
				\item Strategic changes (spin-offs, divestitures)
				\item Governance reforms
			\end{enumerate}
			
			\textbf{Impact on Corporate Governance:}
			\begin{itemize}
				\item Increased accountability
				\item Improved transparency
				\item Enhanced shareholder value
				\item Potential short-term focus
			\end{itemize}
			
			\textbf{Value Creation vs. Value Destruction:}
			\begin{itemize}
				\item Activism can create value by correcting misallocation
				\item However, excessive activism can destroy long-term value
				\item Evidence suggests net positive impact
				\item Companies should engage proactively
			\end{itemize}
		\end{block}
	\end{frame}
	
	%===========================================================================
	% QUESTION 11 - GOVERNANCE REFORM
	%===========================================================================
	\begin{frame}{Question 11: Governance Reform Proposal}
		
		\begin{block}{Comprehensive Governance Reform Proposal}
			
			\subquestion{1. Board Composition}
			\begin{itemize}
				\item Increase independent directors from 60\% to 75\%
				\item Appoint a lead independent director
				\item Implement board evaluations
				\item Add expertise in technology and cybersecurity
			\end{itemize}
			
			\subquestion{2. Committee Structure}
			\begin{itemize}
				\item Strengthen Audit Committee oversight
				\item Create a separate Risk Committee
				\item Enhance Nominating Committee processes
				\item Regular committee evaluations
			\end{itemize}
			
			\subquestion{3. Compensation}
			\begin{itemize}
				\item Adopt balanced scorecard approach
				\item Implement PSUs with ROIC targets
				\item Add clawback provisions
				\item Increase disclosure
			\end{itemize}
			
			\subquestion{4. Shareholder Rights}
			\begin{itemize}
				\item Adopt proxy access
				\item Implement majority voting
				\item Eliminate staggered boards
				\item Enhance transparency
			\end{itemize}
			
			\subquestion{5. Stakeholder Considerations}
			\begin{itemize}
				\item Establish stakeholder advisory council
				\item Integrate ESG into strategy
				\item Enhance sustainability reporting
				\item Engage with all stakeholders
			\end{itemize}
		\end{block}
	\end{frame}
	
	%===========================================================================
	% QUESTION 12 - ESG
	%===========================================================================
	\begin{frame}{Question 12: Stakeholder Value and ESG Integration}
		\begin{block}{Case Context}
			\textbf{Aethelred Industries} is facing pressure from stakeholders to integrate Environmental, Social, and Governance (ESG) factors into its financial strategy.
		\end{block}
		
		\question{Part (a): ESG and Financial Performance}
		
		\textbf{The Shareholder Primacy vs. Stakeholder Theory Debate:}
		
		\subquestion{1. Shareholder Primacy (Friedman, 1970)}
		\begin{itemize}
			\item The only social responsibility of business is to increase profits
			\item ESG activities should only be undertaken if they increase profits
			\item CSR is a distraction from core business
			\item Resources should be allocated to maximize shareholder value
		\end{itemize}
		
		\subquestion{2. Stakeholder Theory (Freeman, 1984)}
		\begin{itemize}
			\item Firms should create value for all stakeholders
			\item This includes employees, customers, communities, environment
			\item Long-term value creation requires stakeholder engagement
			\item Sustainability is essential for long-term viability
		\end{itemize}
	\end{frame}
	
	%===========================================================================
	% QUESTION 12 - ESG EVIDENCE
	%===========================================================================
	\begin{frame}{Question 12: ESG Evidence}
		
		\begin{block}{Empirical Evidence on ESG and Financial Performance}
			
			\textbf{1. Does ESG Create or Destroy Value?}
			\begin{itemize}
				\item \highlight{Positive Relationship:} Studies show positive correlation between ESG and financial performance
				\item \highlight{Causality:} The direction of causality is debated
				\item \highlight{Industry Differences:} ESG impact varies by industry
				\item \highlight{Materiality:} Only material ESG factors affect performance
			\end{itemize}
			
			\textbf{2. Cost of Capital Implications}
			\begin{itemize}
				\item \highlight{Lower Cost of Capital:} High ESG firms have lower cost of equity and debt
				\item \highlight{Better Access:} ESG firms have better access to capital
				\item \highlight{Lower Risk:} ESG factors reduce systematic risk
				\item \highlight{Investor Demand:} Growing demand for ESG investments
			\end{itemize}
			
			\textbf{3. Risk Reduction Benefits}
			\begin{itemize}
				\item \highlight{Operational Risk:} ESG reduces operational risk
				\item \highlight{Reputational Risk:} ESG reduces reputational risk
				\item \highlight{Regulatory Risk:} ESG reduces regulatory risk
				\item \highlight{Systemic Risk:} ESG reduces systemic risk
			\end{itemize}
		\end{block}
	\end{frame}
	
	%===========================================================================
	% QUESTION 12 - ESG STRATEGIES
	%===========================================================================
	\begin{frame}{Question 12: ESG Integration Strategies}
		\question{Part (b): ESG Integration Strategies}
		
		\begin{block}{Proposed ESG Integration Framework}
			
			\subquestion{1. Environmental Strategy}
			\begin{itemize}
				\item \textbf{Carbon Footprint Reduction:}
				\begin{itemize}
					\item Reduce emissions by 30\% by 2030
					\item Invest in renewable energy
					\item Improve energy efficiency
					\item Adopt green technologies
				\end{itemize}
				\item \textbf{Sustainable Supply Chains:}
				\begin{itemize}
					\item Supplier ESG screening
					\item Sustainable sourcing
					\item Circular economy practices
					\item Waste reduction
				\end{itemize}
			\end{itemize}
			
			\subquestion{2. Social Strategy}
			\begin{itemize}
				\item \textbf{Employee Welfare:}
				\begin{itemize}
					\item Competitive compensation
					\item Health and safety
					\item Work-life balance
					\item Career development
				\end{itemize}
				\item \textbf{Community Engagement:}
				\begin{itemize}
					\item Local community investment
					\item Education and training
					\item Philanthropy
				\end{itemize}
			\end{itemize}
		\end{block}
	\end{frame}
	
	%===========================================================================
	% QUESTION 12 - GOVERNANCE STRATEGY
	%===========================================================================
	\begin{frame}{Question 12: Governance Strategy}
		
		\subquestion{3. Governance Strategy}
		\begin{itemize}
			\item \textbf{Board Oversight:}
			\begin{itemize}
				\item ESG committee
				\item Independent ESG expertise
				\item ESG risk oversight
				\item ESG performance monitoring
			\end{itemize}
			\item \textbf{Transparency and Reporting:}
			\begin{itemize}
				\item Sustainability reporting
				\item Stakeholder engagement
				\item ESG disclosure
				\item Assurance of ESG data
			\end{itemize}
			\item \textbf{Ethical Conduct:}
			\begin{itemize}
				\item Code of conduct
				\item Whistleblower protection
				\item Anti-corruption policies
				\item Ethical business practices
			\end{itemize}
		\end{itemize}
		
		\begin{block}{Implementation Roadmap}
			\begin{enumerate}
				\item \textbf{Year 1:} Assessment and baseline
				\item \textbf{Year 2:} Implementation and pilots
				\item \textbf{Year 3:} Full integration and reporting
				\item \textbf{Year 4:} Continuous improvement
			\end{enumerate}
		\end{block}
	\end{frame}
	
	%===========================================================================
	% CONCLUSION
	%===========================================================================
	\begin{frame}{Conclusion: Summary of Key Insights}
		\begin{block}{Key Takeaways from the Synoptic Examination}
			
			\begin{enumerate}
				\item \textbf{Asset Pricing Models:}
				\begin{itemize}
					\item CAPM provides a foundation but is limited
					\item Multifactor models capture additional systematic risks
					\item Real options add value to project evaluation
					\item The choice of model affects capital budgeting decisions
				\end{itemize}
				
				\item \textbf{Market Efficiency:}
				\begin{itemize}
					\item EMH is a useful benchmark but incomplete
					\item AMH provides a more realistic framework
					\item Behavioral factors are permanent features of markets
					\item Understanding market dynamics is essential for success
				\end{itemize}
				
				\item \textbf{Capital Structure:}
				\begin{itemize}
					\item MM propositions provide a foundation
					\item Trade-off theory explains optimal leverage
					\item Pecking order theory explains financing choices
					\item Capital structure decisions are context-dependent
				\end{itemize}
				
				\item \textbf{Corporate Governance:}
				\begin{itemize}
					\item ERM provides a comprehensive risk management framework
					\item Governance mechanisms mitigate agency problems
					\item Compensation structures should align with long-term value
					\item Shareholder engagement creates accountability
				\end{itemize}
				
				\item \textbf{ESG Integration:}
				\begin{itemize}
					\item ESG factors affect financial performance
					\item Integration creates long-term value
					\item Stakeholder engagement is essential
					\item ESG is becoming a competitive necessity
				\end{itemize}
				
				\item \textbf{Neuroeconomics:}
				\begin{itemize}
					\item Decision-making is influenced by neurological factors
					\item Cognitive biases are systematic and predictable
					\item Understanding biases improves decision-making
					\item Organizations should design systems that counteract biases
				\end{itemize}
			\end{enumerate}
		\end{block}
	\end{frame}
	
	%===========================================================================
	% FINAL SLIDE
	%===========================================================================
	\begin{frame}{Final Remarks}
		\begin{center}
			\Huge\textcolor{navyblue}{\textbf{Thank You}}
			\vspace{0.5cm}
			
			\Large\textit{``The modern finance professional must be both a quantitative analyst and a behavioral psychologist, understanding the numbers while accounting for the human factors that drive markets and organizations.''}
			\vspace{0.5cm}
			
			\large\textbf{Shivgan Joshi}
			\vspace{0.2cm}
		\end{center}
	\end{frame}
	
	%===========================================================================
	% REFERENCES
	%===========================================================================
	\begin{frame}{References}
		\scriptsize
		\begin{thebibliography}{99}
			
			\bibitem{beasley2025} Beasley, M. S., \& AICPA. (2025). \textit{2025 The state of risk oversight: An overview of enterprise risk management practices} (16th ed.). Enterprise Risk Management Initiative, NC State University.
			
			\bibitem{bis2025} Bank for International Settlements. (2025, June 24). \textit{Annual economic report 2025: Chapter III – The next-generation monetary and financial system}.
			
			\bibitem{fama1970} Fama, E. F. (1970). Efficient capital markets: A review of theory and empirical work. \textit{The Journal of Finance, 25}(2), 383–417.
			
			\bibitem{fama1993} Fama, E. F., \& French, K. R. (1993). Common risk factors in the returns on stocks and bonds. \textit{Journal of Financial Economics, 33}(1), 3–56.
			
			\bibitem{fama2015} Fama, E. F., \& French, K. R. (2015). A five-factor asset pricing model. \textit{Journal of Financial Economics, 116}(1), 1–22.
			
			\bibitem{kahneman1979} Kahneman, D., \& Tversky, A. (1979). Prospect theory: An analysis of decision under risk. \textit{Econometrica, 47}(2), 263–291.
			
			\bibitem{kuhnen2025} Kuhnen, C. M. (2025). Financial decision-making across the lifespan: Insights from neuroeconomics. \textit{Annual Review of Financial Economics, 17}, 395–410.
			
			\bibitem{lo2004} Lo, A. W. (2004). The adaptive markets hypothesis: Market efficiency from an evolutionary perspective. \textit{The Journal of Portfolio Management, 30}(5), 15–29.
			
			\bibitem{openstax2022} OpenStax. (2022). \textit{Principles of finance}. Rice University.
			
			\bibitem{sharpe1964} Sharpe, W. F. (1964). Capital asset prices: A theory of market equilibrium under conditions of risk. \textit{The Journal of Finance, 19}(3), 425–442.
			
		\end{thebibliography}
	\end{frame}
	
	%===========================================================================
	% END DOCUMENT
	%===========================================================================
\end{document}